
\section{拉格朗日相交中的应用}
\label{sec:6-5}

本节给出弗洛尔上同调在拉格朗日相交问题中的应用。
一方面，是估计 $L$ 与 $\psi(L)$ 的交点数，
其中 $\psi$ 为哈密顿微分同胚；
这将在第 6.5.1、6.5.2 小节讨论。
另一方面，是与哈密顿同痕的霍费尔距离有关的问题，
将在第 6.5.3--6.5.5 小节讨论。

\subsection{定理 H 的证明}
\label{subsec:6-5-1}

现在证明定理 H。
假设
$i_*:H_*(L;\mathbb Q)\to H_*(M;\mathbb Q)$
为单射。
由推论~\ref{cor:3-8-43}，
$L$ 经 bulk 形变后无障碍。
因此可定义满足定理~\ref{thm:D} 的谱序列。
再由 \textup{(D.3)}，
该谱序列在 $E_2$ 页退化。
从而\translatornote{源书实体页 394
在第一个 $HF$ 的秩前
另有一个无指标求和号；
这里取整个 $HF$ 的总秩，
见疑误 A-0391。}
\[
\#(L\cap\psi(L))
\geq
\operatorname{rank}_{\Lambda_{\mathrm{nov}}}
HF\bigl((L,(\boldsymbol b,b)),\psi;\Lambda_{\mathrm{nov}}\bigr)
=\sum_k\operatorname{rank}H^k(L;\mathbb Q).
\]
这里的
$HF((L,(\boldsymbol b,b)),\psi;\Lambda_{\mathrm{nov}})$
如式~\eqref{eq:6-1-8} 所定义。
\hfill$\square$
